\documentclass[11pt]{article}
\usepackage[margin=1in]{geometry}
\usepackage{amsmath,amssymb}
\usepackage{hyperref}

\title{VGB Density-Ratio Decoding: Derivations}
\author{}
\date{}

\newcommand{\pb}{\pi_{\mathrm{big}}}
\newcommand{\ps}{\pi_{\mathrm{small}}}
\newcommand{\Y}{\mathcal{Y}}
\newcommand{\C}{\mathcal{C}}
\newcommand{\R}{\mathcal{R}}

\begin{document}
\maketitle

\noindent\textbf{Reference paper:} \href{https://arxiv.org/pdf/2606.28301}{https://arxiv.org/pdf/2606.28301}

\section{Partial-State Probability}

For a masked state $z$, define
\begin{equation}
\pi(z\mid x)
=
\sum_{y\in\C(z)}\pi(y\mid x),
\end{equation}
where $\C(z)$ is the set of fully revealed sequences compatible with $z$.
Thus $\pi(z\mid x)$ is a partial-state marginal of the full sequence-level joint distribution.

For the fully masked root $z_0$,
\begin{equation}
\C(z_0)=\Y,
\end{equation}
so every normalized model satisfies
\begin{equation}
\pb(z_0\mid x)=\ps(z_0\mid x)=1.
\end{equation}
Hence
\begin{equation}
V(z_0)=1.
\end{equation}

\section{Density-Ratio Value}

Take the small model as the VGB reference model and define the terminal reward
\begin{equation}
\tau(y)
=
\frac{\pb(y\mid x)}{\ps(y\mid x)}.
\end{equation}

Then the reward-tilted target becomes
\begin{align}
\pi^\star(y\mid x)
&\propto
\ps(y\mid x)\tau(y)\\
&=
\pb(y\mid x),
\end{align}
so the target leaf distribution is the big-model distribution.

The exact VGB value is
\begin{align}
V(z)
&=
\mathbb{E}_{Y\sim\ps}
\left[
\tau(Y)\mid Y\in\C(z)
\right]\\
&=
\frac{\pb(z\mid x)}{\ps(z\mid x)}.
\end{align}

\section{Neighboring-State Ratio}

Let $c=z^{i\leftarrow a}$ be obtained by revealing position $i$ as token $a$.
For a coherent joint distribution,
\begin{equation}
\pi(c\mid x)
=
\pi(z\mid x)\pi(Y_i=a\mid x,z).
\end{equation}
Therefore
\begin{align}
\frac{V(c)}{V(z)}
&=
\frac{\pb(c)/\pb(z)}{\ps(c)/\ps(z)}\\
&=
\frac{\pb(Y_i=a\mid x,z)}{\ps(Y_i=a\mid x,z)}.
\end{align}
Thus
\begin{equation}
\boxed{
V(c)
=
V(z)
\frac{\pb(a\mid x,z)}{\ps(a\mid x,z)}.
}
\end{equation}

If $p=z^{-i}$ and $z_i=a$, then equivalently
\begin{equation}
\boxed{
\frac{V(p)}{V(z)}
=
\frac{\ps(a\mid x,p)}{\pb(a\mid x,p)}.
}
\end{equation}

\section{Pathwise Recursion and Caching}

Along a purely forward path
\begin{equation}
z_0\to z_1\to\cdots\to z_k,
\end{equation}
where $z_t$ reveals token $a_t$ from parent $z_{t-1}$,
\begin{equation}
V(z_k)
=
\prod_{t=1}^{k}
\frac{\pb(a_t\mid x,z_{t-1})}{\ps(a_t\mid x,z_{t-1})}.
\end{equation}
In log space,
\begin{equation}
\log V(z_k)
=
\sum_{t=1}^{k}
\left[
\log\pb(a_t\mid x,z_{t-1})
-
\log\ps(a_t\mid x,z_{t-1})
\right].
\end{equation}

A backward move through edge $p\leftrightarrow z$ contributes the reciprocal ratio.
However, an arbitrary remask may create a parent state that was not on the original forward path. Therefore a cached ratio can only be reused when the exact parent context is the same.

Example:
\begin{equation}
[\_,\_,\_,\_]
\to
[A,\_,\_,\_]
\to
[A,B,\_,\_]
\to
[A,B,C,\_].
\end{equation}
If $C$ is remasked, the parent is $[A,B,\_,\_]$, which is the original parent of $C$, so the original edge ratio can be reused.
If $B$ is remasked, the parent becomes $[A,\_,C,\_]$, which is a different context, so a fresh conditional ratio is required.

\section{Forward Transition Simplification}

For a non-leaf state $z$ with $k$ revealed tokens and a child $c=z^{i\leftarrow a}$, Definition 4.2 gives
\begin{equation}
w_{\mathrm{fwd}}(z\to c)
=
 s_{k,1}
 V(z)^\lambda
 \ps(a\mid x,z)
 V(c).
\end{equation}

Using
\begin{equation}
V(c)
=
V(z)
\frac{\pb(a\mid x,z)}{\ps(a\mid x,z)},
\end{equation}
we obtain
\begin{equation}
\boxed{
w_{\mathrm{fwd}}(z\to c)
=
 s_{k,1}
 V(z)^{\lambda+1}
 \pb(a\mid x,z).
}
\end{equation}

\section{Bellman Identity for a Fixed Coordinate}

Fix any masked coordinate $i$.
There is no randomness over the choice of $i$ in the Bellman identity.
The random variable is the eventual completion $Y$ and, when conditioning on position $i$, the token $Y_i$.
By the law of total expectation,
\begin{equation}
\boxed{
V(z)
=
\sum_a
\ps(Y_i=a\mid x,z)
V(z^{i\leftarrow a})
}
\end{equation}
for every fixed unrevealed coordinate $i$.

Hence the total forward weight associated with one fixed coordinate is
\begin{align}
F_i(z)
&=
\sum_a
 s_{k,1}
 V(z)^\lambda
 \ps(a\mid x,z)
 V(z^{i\leftarrow a})\\
&=
 s_{k,1}V(z)^{\lambda+1}.
\end{align}

Since there are exactly $n-k$ masked coordinates,
\begin{equation}
\boxed{
F(z)
=
(n-k)s_{k,1}V(z)^{\lambda+1}.
}
\end{equation}

This is not an expectation over $i$; it is the sum of $n-k$ coordinate-specific totals, each equal to the same quantity by the Bellman identity.

\section{Backward Transition Simplification}

For parent $p_i=z^{-i}$, Definition 4.2 gives
\begin{equation}
w_{\mathrm{bwd}}(z\to p_i)
=
 s_{k-1,1}
 V(z)
 V(p_i)^\lambda.
\end{equation}

Define
\begin{equation}
\rho_i
:=
\frac{V(p_i)}{V(z)}.
\end{equation}
Then
\begin{equation}
w_{\mathrm{bwd}}(z\to p_i)
=
 s_{k-1,1}
 V(z)^{\lambda+1}
 \rho_i^\lambda.
\end{equation}

Under the density-ratio value,
\begin{equation}
\boxed{
\rho_i
=
\frac{\ps(Y_i=z_i\mid x,p_i)}{\pb(Y_i=z_i\mid x,p_i)}.
}
\end{equation}

Therefore
\begin{equation}
\boxed{
B(z)
=
 s_{k-1,1}
 V(z)^{\lambda+1}
 \sum_{i\in\R(z)}\rho_i^\lambda.
}
\end{equation}

\section{Derivation of \texorpdfstring{$W(z)$}{W(z)}}

By definition,
\begin{equation}
W(z)=F(z)+B(z).
\end{equation}
Substituting the preceding expressions,
\begin{equation}
\boxed{
W(z)
=
V(z)^{\lambda+1}
\left[
(n-k)s_{k,1}
+
 s_{k-1,1}
\sum_{i\in\R(z)}\rho_i^\lambda
\right].
}
\end{equation}

Define the reduced normalizer
\begin{equation}
\boxed{
\widetilde W(z)
=
(n-k)s_{k,1}
+
 s_{k-1,1}
\sum_{i\in\R(z)}\rho_i^\lambda.
}
\end{equation}

Then
\begin{align}
P(\text{forward}\mid z)
&=
\frac{(n-k)s_{k,1}}{\widetilde W(z)},\\
P(\text{backward}\mid z)
&=
\frac{s_{k-1,1}\sum_i\rho_i^\lambda}{\widetilde W(z)}.
\end{align}

The absolute factor $V(z)^{\lambda+1}$ cancels. Therefore the sampler does not need to evaluate or recursively maintain $V(z)$ itself.

\section{Conditional Distributions Within Each Direction}

For a forward move, one fixed coordinate has total mass
\begin{equation}
s_{k,1}V(z)^{\lambda+1},
\end{equation}
independent of which masked coordinate is chosen. Hence
\begin{equation}
P(i\mid \text{forward},z)=\frac{1}{n-k}.
\end{equation}

Given coordinate $i$,
\begin{equation}
P(a\mid i,\text{forward},z)
=
\pb(Y_i=a\mid x,z).
\end{equation}

For a backward move,
\begin{equation}
P(i\mid \text{backward},z)
=
\frac{\rho_i^\lambda}{\sum_j\rho_j^\lambda}.
\end{equation}

\section{Implementation Consequence}

The theoretically awkward partial-state joint ratio
\begin{equation}
V(z)=\frac{\pb(z)}{\ps(z)}
\end{equation}
never needs to be evaluated explicitly for normalized non-leaf transitions.
Only local leave-one-out ratios
\begin{equation}
\frac{\ps(z_i\mid z^{-i})}{\pb(z_i\mid z^{-i})}
\end{equation}
are required for backward weighting, together with the big-model conditional distribution for the selected forward coordinate.

A practical implementation can therefore construct a batch consisting of the current state and all or a shortlisted set of remasked parents, and evaluate that batch once with each model.

\section{Model-Consistency Caveat}

All equalities above are exact when the conditional probabilities produced by each model are compatible with a coherent underlying sequence-level joint distribution. A practical masked language model may only approximately satisfy this compatibility across arbitrary mask patterns and reveal orders. In that case, the local density-ratio rule remains implementable but should be viewed as an approximation to the exact state-density-ratio VGB construction.

\end{document}
